\documentclass{article}
\usepackage[portuguese]{babel}
\usepackage{pgfplots}

\title{Trabalho Prático - Confiabilidade e Segurança de Software}

\author{Pedro Arthur Pamplona Hartmann}
\date{}

\begin{document}
\maketitle
\begin{abstract}
O objetivo deste trabalho é atacar cenários onde seria utilizado o protocolo Diffie-Hellman.
O ataque foi feito usando técnicas conhecidas para resolver o problema do logaritmo discreto, o Pohlig-Hellman e Pollard rho para logaritmos.
O problema do logaritmo discreto é definido da seguinte forma.
Temos elementos $\alpha,\beta,n$ tal que $\beta \equiv \alpha^\gamma \pmod n$ e o objetivo é achar $\gamma$.
\end{abstract}

\section{Ataques utilizados}
\subsection{Pohlig-Hellman}
Pohlig-Hellman é um algoritmo que divide o problema do logaritmo discreto original em problemas menores e usa o Teorema Chinês do Resto para combinar os resultados.
Cada sub-problema foi resolvido usando Pollard rho.

\subsection{Pollard rho para logaritmos}
Para achar o logaritmo, o algoritmo tenta achar $a,b,A,B$ tal que $\alpha^a\beta^b=\alpha^A\beta^B$, usando algum algoritmo de detecção de ciclos.
Foi definido uma função $f:x_i \mapsto x_{i+1}$ que define uma sequência pseudo-aleatória onde a invariante é $x_i=\alpha^{a_i}\beta^{b_i}$.
Para achar uma colisão foi utilizado uma técnica proposta por van Oorschot e Wiener \cite{oorschot+wiener}.
A implementação utilizada da técnica usa uma tabela hash para guardar alguns elementos da sequência dado por uma propriedade facilmente testável.
Quando um $x_i$ com essa propriedade é encontrado o algoritmo verifica se esse número está na tabela.
Se o estiver na tabela, há uma possível colisão.
Se não estiver insere.
Se houver uma colisão, temos uma equação onde o $\gamma$ pode ser recuperado.

\subsection{Complexidade}
A complexidade dos algoritmos combinados é de $\mathcal{O}(\sqrt p)$ onde $p$ é maior fator primo de $n$.

\section{Gráfico de execução}
\begin{tikzpicture}
\begin{semilogyaxis} [
	xlabel={Bits em $p$},
	ylabel={Tempo [Segundos]},
	xmin=32, xmax=1024,
	ymin=0, ymax=3600,
]

\addplot[
	color=blue,
	mark=o,
	]
	coordinates{
		(32, 0.002)(40, 0.0005)(48, 0.001)(64, 0.008)(80,0.02)(96,0.04)(112,0.1)(128,0.49)
	};

\addplot[
	color=red,
	mark=o,
	]
	coordinates{
		(32, 0.004)(40, 0.03)(48, 0.3)(64, 89.1)
	};

\end{semilogyaxis}
\end{tikzpicture}

\section{Infraestrutura Computacional}
\begin{itemize}
	\item Modelo da CPU: Intel i7-1255U.
	\item 12 núcleos lógicos no total.
	\item P-cores: 2 núcleos físicos e 4 núcleos lógicos.
	\item E-cores: 8 núcleos físicos.
	\item Memória RAM: 16 GiB.
	\item Sistema operacional: Linux 7.0.10
\end{itemize}

\nocite{hac}
\bibliographystyle{plain}
\bibliography{citations}

\end{document}
